%% kapa: exam=1 round=34 year=2023 period=1 form=A subject=civil_law subject_ko=민법 seq=38 number=38
%% source: Q-Net 감정평가사 기출문제 (cst00309) artlSeq=5211789 file=raw/1008/5211789/2023년도 제34회 감정평가사 1차시험 1교시.pdf
%% conversion: kordoc 4.24.51 -> tools/gen1.py (auto); review=proofread-sonnet; flags=none
\begin{question}{38}
저당권에 관한 설명으로 옳지 않은 것은? (다툼이 있으면 판례에 따름)
\begin{choices}
  \item 채권자와 제3자가 불가분적 채권자의 관계에 있다고 볼 수 있는 경우에는 그 제3자 명의의 저당권등기도 유효하다.
  \item 근저당권설정자가 적법하게 기본계약을 해지하면 피담보채권은 확정된다.
  \item 무효인 저당권등기의 유용은 그 유용의 합의 전에 등기상 이해관계가 있는 제3자가 없어야 한다.
  \item 저당부동산의 제3취득자는 부동산의 개량을 위해 지출한 유익비를 그 부동산의 경매대가에서 우선 변제받을 수 없다.
  \item 저당권자가 저당부동산을 압류한 이후에는 저당권설정자의 저당부동산에 관한 차임채권에도 저당권의 효력이 미친다.
\end{choices}
\end{question}
\answer{4}
